% !TEX root = ../mainfile.tex
% !TEX encoding = UTF-8 Unicode
% !TEX spellcheck = en_US

\chapter{Conclusion}
\label{ch:conclusion}

This thesis asked whether the U.S.\ dollar's safe-haven response depends on the type of
shock that triggers a crisis, even when the United States is implicated in the shock itself.
Two recent episodes provided an unusually clean test: the 2025 Liberation Day tariffs, a
purely U.S.\ economic-policy shock, and the 2026 Iran/Hormuz crisis, a geopolitical
oil-supply shock, with the 2022 Russian invasion of Ukraine as an external control.

The evidence gives a clear answer. The two shock types produced a mirror image. The dollar
depreciated around the tariff shock---roughly four percent on the broad index within twenty
trading days---while it appreciated around the Hormuz crisis and around Ukraine, and a formal
difference test confirms the tariff response is significantly distinct from both security
episodes, which are themselves statistically indistinguishable from one another. The dollar's
haven status is therefore conditional on the \emph{nature} of the shock, not on its origin:
it survives a U.S.-originated security shock but inverts for a U.S.-originated economic-policy
shock. A complementary regression confirms the breakdown statistically: the dollar's
safe-haven sensitivity to risk---significantly positive on normal days---collapses to zero in
the tariff window, while holding through the Hormuz and Ukraine episodes. Two further results
round out the picture:
gold acts as a hedge against the dollar rather than as a universal safe haven, and the oil
channel sorts currencies only for the supply-driven shock.

The thesis contributes the first systematic side-by-side comparison of the Liberation Day and
Iran/Hormuz episodes for currency markets, carries a gold and oil benchmark through the same
design to connect the safe-haven and commodity-currency literatures, and extends standard
cross-sectional currency tools to mid-2026 on a consistent WMR dataset.

Several extensions suggest themselves. As the geopolitical situation continues to evolve, the
sample can be extended so that the most recent episodes acquire full post-event windows.
Richer forward and basis data would allow covered-interest-parity deviations to be brought
into the analysis, and intraday data would tighten the identification around the announcement
windows. A dynamic conditional-correlation (DCC-GARCH) model of the dollar--VIX relationship
would test the safe-haven-breakdown hypothesis directly and is the natural next layer.
Finally, a formal model of shock-dependent safe-haven demand---in which the currency of the
shock's originator loses its haven premium for policy shocks but not for security
shocks---would give the empirical pattern documented here a theoretical foundation.
